\begin{pastexam}[label={px:ntu-107-ie-ds-10}]{臺大工工所 107 學年度 第 10 題}
（是非；命題正確填 (A)，錯誤填 (B)；答對得 4 分，答錯倒扣 4 分，至多扣到該大題零分）

In an AVL tree, every subtree is also an AVL tree.
\end{pastexam}

\begin{pastsoln}
答案是 (A)，正確。

AVL 樹的條件是兩件事（資結書第~10.2~節）：它是 BST，而且\textbf{每個節點}的左右子樹高度差不超過
$1$。 取任何一個節點 $x$ 為根的子樹：
\begin{itemize}
\item BST 性質：子樹裡每個節點的左子樹都比它小、右子樹都比它大，這在整棵樹成立，在子樹裡當然成立。
\item 平衡條件：子樹裡每個節點的左右子樹，和它在整棵樹裡的左右子樹是同一棵，高度不變，高度差仍
然 $\le 1$。
\end{itemize}
兩個條件都只看「節點和它底下的東西」，和上面的祖先無關，所以每棵子樹都是 AVL 樹。 這一點和下面
第 13 題的紅黑樹不同：紅黑樹多了一條「根是黑色」，那條會被子樹破壞。
\end{pastsoln}
